\documentclass[a4paper]{article}
% Español (México) con XeLaTeX: fontspec para fuentes Unicode, polyglossia
% para la división silábica y los nombres del documento en la variante mexicana.
\usepackage{fontspec}
\usepackage{polyglossia}
\setdefaultlanguage[variant=mexican]{spanish}
% Los nombres chinos van como «pinyin (original)»: xeCJK da a los caracteres
% Han su propia fuente; sin ella XeLaTeX los omite con un aviso.
% FandolSong no cubre algunos caracteres de nombres (U+73FA); WenQuanYi Zen Hei sí.
\usepackage[AutoFallBack=true]{xeCJK}
\setCJKmainfont{FandolSong-Regular.otf}
\setCJKfallbackfamilyfont{\CJKrmdefault}{WenQuanYi Zen Hei}
% polyglossia no traduce los nombres de listings.
\AtBeginDocument{\providecommand{\lstlistingname}{}\renewcommand{\lstlistingname}{Listado}%
  \providecommand{\lstlistlistingname}{}\renewcommand{\lstlistlistingname}{Índice de listados}}
\usepackage{amsmath, amssymb}
\usepackage{graphicx}
\usepackage[margin=2.5cm]{geometry}
\usepackage[most]{tcolorbox}
\usepackage{etoolbox}
\usepackage{listings}
\usepackage{hyperref}
\usepackage{booktabs}
\usepackage{float}
\newtcolorbox{knowledgebox}[1]{enhanced, colback=blue!5!white, colframe=blue!75!black, colbacktitle=blue!75!black, coltitle=white, fonttitle=\bfseries, title={#1}, attach boxed title to top left={yshift=-2mm, xshift=2mm}, boxrule=1pt, sharp corners}
\newtcolorbox{importantbox}[1]{enhanced, colback=yellow!10!white, colframe=yellow!80!black, colbacktitle=yellow!80!black, coltitle=black, fonttitle=\bfseries, title={#1}, sharp corners}
\newtcolorbox{warningbox}[1]{enhanced, colback=red!5!white, colframe=red!75!black, colbacktitle=red!75!black, coltitle=white, fonttitle=\bfseries, title={#1}, sharp corners}
\lstset{language=Python, basicstyle=\ttfamily\small, keywordstyle=\color{blue}, stringstyle=\color{red!60!black}, commentstyle=\color{green!60!black}, breaklines=true, frame=single, numbers=left, numberstyle=\tiny\color{gray}, captionpos=b, extendedchars=true}
\newcommand{\notetitle}{Práctica del ciclo agéntico de veRL}
\newcommand{\noteauthors}{Elaboradas a partir de materiales públicos del curso}
\newcommand{\notedate}{2025}
\newcommand{\videochannel}{Wudaokou Nashi (五道口纳什)}
\newcommand{\videopublishdate}{2025}
\newcommand{\videoduration}{aproximadamente 20 minutos}
\newcommand{\videourl}{}
\newcommand{\videocoverpath}{cover.jpg}
\newcommand{\repourl}{https://github.com/hqhq1025/ai-course-notes}

% Footer with repo info
\usepackage{fancyhdr}
\pagestyle{fancy}
\fancyhf{}
\fancyfoot[C]{\thepage}
\fancyfoot[R]{\footnotesize\color{gray}\href{\repourl}{AI Course Notes}}
\renewcommand{\headrulewidth}{0pt}
\renewcommand{\footrulewidth}{0pt}

\begin{document}
\begin{titlepage}\centering
{\Large Notas del curso\par}\vspace{1.2cm}{\huge\bfseries \notetitle\par}\vspace{0.8cm}{\large \noteauthors\par}\vspace{0.3cm}{\large \notedate\par}\vspace{1.2cm}
\ifdefempty{\videocoverpath}{}{\includegraphics[width=0.82\textwidth,height=0.45\textheight,keepaspectratio]{\videocoverpath}\par}
\vfill
\begin{tcolorbox}[width=0.9\textwidth, colback=black!2!white, colframe=black!60, sharp corners]
\textbf{Autor o canal del video}: \videochannel\par\textbf{Fecha de publicación}: \videopublishdate\par\textbf{Duración del video}: \videoduration\par
\ifdefempty{\videourl}{}{\textbf{Enlace del video}: \href{\videourl}{\nolinkurl{\videourl}}\par}\par
\textbf{Página del proyecto}: \href{\repourl}{\nolinkurl{\repourl}}\par
\end{tcolorbox}
\end{titlepage}
\tableofcontents\newpage

\section{Introducción}
Esta sesión entra de lleno en el núcleo del entrenamiento de Agentic RL: la implementación del Agentic Loop en veRL.

\begin{importantbox}{El concepto del Agentic Loop}
El agente interactúa con el entorno mediante un ciclo de varias rondas de think-act-observe. En el entrenamiento de RL, el modelo necesita, en cada paso:
\begin{enumerate}
  \item Generar (Generate): producir texto o una tool call
  \item Llamar a la herramienta (Tool Call): ejecutar código, consultar una API, etc.
  \item Obtener retroalimentación (Observe): recibir el resultado que devuelve la herramienta
  \item Continuar generando: seguir razonando a partir de la retroalimentación
\end{enumerate}
Este proceso de Multi-turn se controla mediante el archivo de configuración.
\end{importantbox}

\section{Arquitectura del Agentic Loop de veRL}
Archivo de código central: \texttt{agentloop.py}, que incluye:
\begin{itemize}
  \item \texttt{AgentLoopManager}: administra de manera unificada todos los procesos
  \item \texttt{AgentLoopWorker}: ejecuta la lógica concreta del Agent
  \item \texttt{AgentLoopBase}: clase base abstracta, requiere personalizar el método \texttt{run}
  \item Server Manager: administra el servicio de inferencia
\end{itemize}

\section{Método generate\_sequence}
El método \texttt{generate\_sequence} de \texttt{AgentLoopManager} es el núcleo, ya que unifica la gestión de todos los componentes.

\subsection{Resumen de la sección}
El Agentic Loop de veRL ofrece un marco completo de entrenamiento con RL de múltiples turnos.

\section{Perspectiva del sistema: la división del trabajo entre Manager, Worker y Loop}
Entender el Agentic Loop de veRL no consiste solo en memorizar unos cuantos nombres de clases, sino en ver con claridad cómo separa el razonamiento de múltiples turnos, la invocación de herramientas y la gestión de servicios. El Manager se encarga de la orchestration, el Worker se encarga de la ejecución y LoopBase se encarga de definir la interfaz abstracta; esta división en capas permite que el sistema reutilice la infraestructura entre distintas tareas.

\begin{knowledgebox}{Por qué esta arquitectura conviene al Agentic RL}
\begin{itemize}
  \item Un bucle de múltiples turnos requiere de manera natural una gestión de estado unificada
  \item La invocación de herramientas y el razonamiento del modelo por lo general deben reutilizarse en módulos separados
  \item Durante el entrenamiento hay que atender al mismo tiempo el rollout, el server y la lógica de la tarea
\end{itemize}
\end{knowledgebox}

\subsection{Resumen de la sección}
Lo importante del Agentic Loop no es qué tan compleja sea alguna clase en particular, sino que abstrae la lógica de interacción think-act-observe en un esqueleto de sistema reutilizable.

\newpage
\section{Síntesis y ampliación}
\begin{enumerate}
  \item Agentic Loop = ciclo think-act-observe de Multi-turn
  \item veRL se administra de forma unificada mediante AgentLoopManager
  \item La lógica del Agent se puede personalizar (SingleTool, TwoAgents, etc.)
\end{enumerate}
\end{document}

